\documentclass[11pt]{article}
\usepackage[a4paper,margin=22mm]{geometry}
\usepackage{fontspec}
\setmainfont{DejaVu Serif}
\setsansfont{DejaVu Sans}
\usepackage{microtype}
\usepackage{xcolor}
\usepackage{hyperref}
\usepackage{enumitem}
\usepackage{parskip}
\definecolor{Accent}{HTML}{971D32}
\definecolor{Muted}{HTML}{666666}
\hypersetup{colorlinks=true,urlcolor=Accent,linkcolor=Accent}
\setlist{leftmargin=1.6em,itemsep=0.3em,topsep=0.35em}
\setcounter{tocdepth}{2}
\renewcommand{\contentsname}{Contents}
\newcommand{\sectionrule}{\par\vspace{-0.5em}\noindent\textcolor{Accent}{\rule{\linewidth}{0.6pt}}\par}
\begin{document}

\begin{center}
{\sffamily\bfseries\color{Accent} TOEFL WORD OF THE DAY\par}
\vspace{8mm}
{\fontsize{42}{48}\selectfont\bfseries calamity\par}
\vspace{2mm}
{\Large\sffamily\color{Accent} /kəˈlæməti/\par}
\vspace{2mm}
{\small\sffamily Local audio phrase: calamity\par}
{\small\sffamily Audio file: audios/calamity.mp3\par}
\vspace{2mm}
{\large\sffamily\color{Muted} countable noun\par}
\vspace{5mm}
{\sffamily a terrible and damaging event\par}
\vspace{6mm}
{\large A calamity is a very serious event that causes widespread damage, suffering, or loss.\par}
\vspace{6mm}
\begin{minipage}{0.84\textwidth}
\itshape “Crop failure could become a national calamity if food supplies are already limited.”
\end{minipage}\par
\vspace{10mm}
\textbf{Date:} August 29th 2026\quad\textbullet\quad \textbf{Level:} B1--mid-B2 TOEFL
\vfill
Created and compiled by \textbf{Amir Kooshky}\par
\vspace{2mm}
\href{https://t.me/KooshkyTOEFL}{Telegram Channel}\quad\textbullet\quad
\href{https://www.instagram.com/kooshkytoefl}{Instagram}\quad\textbullet\quad
\href{https://t.me/KooshkyTOEFL\_pv}{Personal Telegram}
\end{center}
\newpage
\tableofcontents
\newpage

\section{Meanings and usage}
\sectionrule
\subsection{Meaning 1: Terrible and damaging event}
\textbf{Part of speech:} countable noun

\textbf{IPA:} /kəˈlæməti/

\textbf{Pronunciation phrase:} calamity

\textbf{Local audio file:} audios/calamity.mp3

\textbf{Register:} formal and somewhat literary; common in serious discussions of disasters, history, war, famine, and large-scale suffering

\textbf{Definition:} a terrible event that causes great damage, suffering, or loss

\textbf{In simpler words:} a terrible disaster

\textbf{Typical pattern:} a major, great, national, or natural calamity

\subsubsection{Natural examples}
\begin{enumerate}
  \item The earthquake was a major calamity that left thousands of people without homes.
  \item Crop failure could become a national calamity if food supplies are already limited.
  \item Better warning systems can help communities prepare for natural calamities.
  \item The destruction of the harvest was a calamity for families who depended entirely on farming.
  \item Officials hoped that early evacuation would prevent the approaching storm from becoming a human calamity.
  \item The collapse of the irrigation system would have been a calamity for the entire region.
  \item Historians disagree about whether the famine was an unavoidable natural calamity or the result of poor government decisions.
  \item Years of careful planning helped the city avoid a far greater calamity when the river flooded.
\end{enumerate}

\subsubsection{Nuance and implication}
\textbf{Central nuance:} Calamity emphasizes the severe suffering, damage, or loss caused by an event rather than simply the fact that something went wrong.

\textbf{Usual implication:} The event is serious enough to affect many people or have major consequences.

\textbf{Compared with minor setback:} A minor setback causes a limited problem or delay. A calamity causes severe damage, suffering, or loss.

\subsubsection{Grammar notes}
\textbf{How to use it:} Calamity is usually a countable noun. Say a calamity for one event and calamities for more than one.

\textbf{What can follow it:} It is often modified by words such as natural, national, major, great, terrible, or widespread.

\textbf{Common preposition:} Use a calamity for someone or something when naming who is badly affected, as in The drought was a calamity for local farmers.

\textbf{Noun use:} The singular normally needs a determiner such as a, the, or this. The plural is calamities.

\subsubsection{Useful patterns}
\textbf{Pattern:} a natural calamity

\textbf{Use:} Use this for a disastrous event caused by natural forces.

\textbf{Example:} The region has experienced several natural calamities, including floods and earthquakes.

\textbf{Pattern:} a national calamity

\textbf{Use:} Use this for a disaster severe enough to affect an entire country or be regarded as a major national tragedy.

\textbf{Example:} The famine developed into a national calamity.

\textbf{Pattern:} a calamity for + person or group

\textbf{Use:} Use this to identify the people who suffer badly because of an event.

\textbf{Example:} The loss of the bridge was a calamity for isolated rural communities.

\textbf{Pattern:} avert or prevent a calamity

\textbf{Use:} Use this when action stops a serious disaster from occurring.

\textbf{Example:} Emergency repairs helped avert a much greater calamity.

\textbf{Pattern:} become or turn into a calamity

\textbf{Use:} Use this when a difficult situation develops into a serious disaster.

\textbf{Example:} Without immediate assistance, the food shortage could become a humanitarian calamity.

\subsubsection{Common wrong usage}
\paragraph{Correction 1}
\textbf{Wrong:} Missing the bus was a terrible calamity for me.

\textbf{Correct:} Missing the bus was an inconvenience for me.

\textbf{Why:} Calamity is normally reserved for events causing very serious damage, suffering, or loss, so it is too strong for an ordinary inconvenience.

\paragraph{Correction 2}
\textbf{Wrong:} The earthquake caused a great calamity disaster.

\textbf{Correct:} The earthquake caused a great calamity.

\textbf{Why:} Calamity already means a serious disaster, so calamity disaster is unnecessary.

\paragraph{Correction 3}
\textbf{Wrong:} The country experienced many calamitys.

\textbf{Correct:} The country experienced many calamities.

\textbf{Why:} The plural of calamity is calamities.

\paragraph{Correction 4}
\textbf{Wrong:} The drought was calamity for local farmers.

\textbf{Correct:} The drought was a calamity for local farmers.

\textbf{Why:} Calamity is countable, so the singular normally needs a determiner such as a.

\section{Word family}
\sectionrule
\subsection{calamitous}
\textbf{Part of speech:} adjective

\textbf{IPA:} /kəˈlæmətəs/

\textbf{Definition:} causing or involving very great damage, suffering, or loss

\textbf{Relationship to the headword:} This adjective describes an event, result, or situation that causes the kind of severe damage associated with a calamity.

\textbf{Grammar pattern:} calamitous + consequences, effects, decline, event, or decision

\textbf{Usage note:} It is formal and much less common in everyday speech than words such as disastrous.

\subsubsection{Examples}
\begin{enumerate}
  \item The war had calamitous effects on the country's economy.
  \item A prolonged failure of the water supply would have calamitous consequences for the city.
  \item The species suffered a calamitous decline after its habitat was destroyed.
\end{enumerate}

\subsection{calamitously}
\textbf{Part of speech:} adverb

\textbf{IPA:} /kəˈlæmətəsli/

\textbf{Definition:} in a way that causes very serious damage, suffering, or loss

\textbf{Relationship to the headword:} This adverb describes an action or event that ends or develops with disastrous consequences.

\textbf{Grammar pattern:} fail, end, or go wrong + calamitously

\textbf{Usage note:} It is uncommon and mainly found in formal writing.

\subsubsection{Examples}
\begin{enumerate}
  \item The policy failed calamitously and created severe shortages.
  \item The expedition ended calamitously after supplies were lost in a storm.
\end{enumerate}

\section{Common collocations}
\sectionrule
\subsection{natural calamity}
\textbf{Applies to:} Terrible and damaging event

\textbf{Meaning:} a severe disaster caused by natural forces

\textbf{Example:} The government created an emergency fund for communities affected by natural calamities.

\subsection{national calamity}
\textbf{Applies to:} Terrible and damaging event

\textbf{Meaning:} a disaster with major consequences for an entire country

\textbf{Example:} The prolonged famine was regarded as a national calamity.

\subsection{major calamity}
\textbf{Applies to:} Terrible and damaging event

\textbf{Meaning:} a very serious disaster

\textbf{Example:} Quick action prevented the industrial accident from becoming a major calamity.

\subsection{great calamity}
\textbf{Applies to:} Terrible and damaging event

\textbf{Meaning:} an event causing severe damage or suffering

\textbf{Example:} The destruction of the city's water system would have been a great calamity.

\subsection{humanitarian calamity}
\textbf{Applies to:} Terrible and damaging event

\textbf{Meaning:} a disaster involving severe human suffering

\textbf{Example:} Aid organizations warned that the conflict could produce a humanitarian calamity.

\subsection{avert a calamity}
\textbf{Applies to:} Terrible and damaging event

\textbf{Meaning:} prevent a serious disaster from happening

\textbf{Example:} Engineers repaired the dam in time to avert a calamity.

\subsection{prevent a calamity}
\textbf{Applies to:} Terrible and damaging event

\textbf{Meaning:} stop a serious disaster from occurring

\textbf{Example:} Early detection of the disease helped prevent a much larger calamity.

\subsection{calamitous consequences}
\textbf{Applies to:} Terrible and damaging event

\textbf{Meaning:} extremely damaging results

\textbf{Example:} A complete failure of the power grid could have calamitous consequences.

\textbf{Usage note:} Calamitous is the adjective form.

\subsection{calamitous effects}
\textbf{Applies to:} Terrible and damaging event

\textbf{Meaning:} extremely harmful effects

\textbf{Example:} The invasion had calamitous effects on agriculture and trade.

\textbf{Usage note:} Calamitous is the adjective form.

\section{Synonyms and antonyms}
\sectionrule
\subsection{Close synonyms}
\subsubsection{disaster}
\textbf{Part of speech:} countable noun

\textbf{Applies to:} Terrible and damaging event

\textbf{Shared meaning:} Both describe events that cause serious damage, suffering, or loss.

\textbf{Difference:} Disaster is the more common and general word. Calamity is more formal and often emphasizes widespread suffering or severe consequences.

\textbf{Example:} The earthquake was one of the worst natural disasters in the country's history.

\subsubsection{catastrophe}
\textbf{Part of speech:} countable noun

\textbf{Applies to:} Terrible and damaging event

\textbf{Shared meaning:} Both refer to very serious disastrous events.

\textbf{Difference:} Catastrophe usually suggests an extremely severe or sudden disaster with enormous consequences. Calamity is also very serious but does not always imply the same extreme scale.

\textbf{Example:} A sudden collapse of the dam would create an environmental catastrophe.

\subsubsection{tragedy}
\textbf{Part of speech:} countable noun

\textbf{Applies to:} Terrible and damaging event

\textbf{Shared meaning:} Both can describe events involving severe suffering and loss.

\textbf{Difference:} Tragedy emphasizes human suffering, death, or deep sadness. Calamity emphasizes severe damage and misfortune and can apply to social, economic, or natural disasters.

\textbf{Example:} The deaths of so many civilians were a terrible tragedy.

\subsubsection{cataclysm}
\textbf{Part of speech:} countable noun

\textbf{Applies to:} Terrible and damaging event

\textbf{Shared meaning:} Both can refer to extremely destructive events.

\textbf{Difference:} Cataclysm is stronger and less common, usually referring to a sudden event that causes enormous destruction or major social change. Calamity is broader.

\textbf{Example:} The eruption was a geological cataclysm that transformed the surrounding landscape.

\section{Words learners may confuse}
\sectionrule
\subsection{casualty}
\textbf{Part of speech:} countable noun

\textbf{Why learners confuse them:} Both frequently appear in discussions of disasters and serious accidents.

\textbf{Key difference:} A calamity is the disastrous event itself. A casualty is a person who is killed or injured in an accident, war, or disaster, or sometimes something lost or damaged because of an event.

\textbf{calamity:} The earthquake was a major calamity for the region.

\textbf{casualty:} The earthquake caused hundreds of casualties.

\subsection{catastrophe}
\textbf{Part of speech:} countable noun

\textbf{Why learners confuse them:} The two words overlap strongly and are often used in similar serious contexts.

\textbf{Key difference:} Both words describe serious disasters, but catastrophe often suggests even greater or more sudden destruction. Calamity is formal and strongly emphasizes serious misfortune and suffering.

\textbf{calamity:} The crop failure became a national calamity.

\textbf{catastrophe:} Scientists warned that complete collapse of the ecosystem would be a global catastrophe.

\subsection{calamitous}
\textbf{Part of speech:} noun versus adjective

\textbf{Why learners confuse them:} They belong to the same word family but have different grammatical roles.

\textbf{Key difference:} Calamity is a noun naming a disastrous event. Calamitous is an adjective describing something that causes or involves disastrous consequences.

\textbf{calamity:} The flood was a calamity for the farming communities.

\textbf{calamitous:} The flood had calamitous effects on agriculture.

\section{Etymology}
\sectionrule
\textbf{Origin:} Calamity comes from Latin calamitas, meaning damage, loss, or disaster.

\textbf{Memory link:} Think of a calamity as an event serious enough to bring major damage, loss, or suffering.

\section{Practice exercises}
\sectionrule
\subsection{Word family choice}
Choose the best member of the word family for each blank.

\begin{enumerate}
  \item The collapse of the irrigation system would have \_\_\_\_\_ consequences for the farming region.
  \item The expedition ended \_\_\_\_\_ after the group lost most of its supplies.
  \item The drought developed into a national \_\_\_\_\_.
\end{enumerate}

\subsection{Correct or incorrect?}
Decide whether each sentence is correct. Correct the inaccurate ones.

\begin{enumerate}
  \item The earthquake was a major calamity for the region.
  \item Forgetting my pencil was a serious calamity.
  \item The country experienced several natural calamities during the century.
  \item The drought was calamity for local farmers.
  \item Emergency measures helped avert a much greater calamity.
\end{enumerate}

\subsection{Collocation practice}
Complete each sentence with a natural collocation.

\begin{enumerate}
  \item Early evacuation helped \_\_\_\_\_ a much greater calamity.
  \item The prolonged famine became a national \_\_\_\_\_.
  \item The storm was a calamity \_\_\_\_\_ communities that depended on the coastal fishing industry.
  \item The policy failure had \_\_\_\_\_ consequences for the country's food supply.
\end{enumerate}

\subsection{Complete answer key}
\subsubsection{Word family choice}
\begin{enumerate}
  \item \textbf{calamitous.} The collapse of the irrigation system would have calamitous consequences for the farming region. \emph{Use the adjective calamitous before consequences.}
  \item \textbf{calamitously.} The expedition ended calamitously after the group lost most of its supplies. \emph{Use the adverb calamitously to describe how the expedition ended.}
  \item \textbf{calamity.} The drought developed into a national calamity. \emph{Use the noun calamity after a national.}
\end{enumerate}

\subsubsection{Correct or incorrect?}
\begin{enumerate}
  \item \textbf{Correct} \emph{Calamity naturally describes an event causing widespread damage and suffering.}
  \item \textbf{Incorrect}. Forgetting my pencil was a minor inconvenience. \emph{Calamity is much too strong for an ordinary small problem.}
  \item \textbf{Correct} \emph{Calamity is countable, and its plural is calamities.}
  \item \textbf{Incorrect}. The drought was a calamity for local farmers. \emph{The singular countable noun calamity normally needs a determiner such as a.}
  \item \textbf{Correct} \emph{Avert a calamity is a natural formal collocation meaning prevent a disaster.}
\end{enumerate}

\subsubsection{Collocation practice}
\begin{enumerate}
  \item \textbf{avert.} Early evacuation helped avert a much greater calamity. \emph{Avert a calamity means prevent a serious disaster.}
  \item \textbf{calamity.} The prolonged famine became a national calamity. \emph{A national calamity is a disaster with severe consequences for an entire country.}
  \item \textbf{for.} The storm was a calamity for communities that depended on the coastal fishing industry. \emph{Use a calamity for someone or something to identify who was badly affected.}
  \item \textbf{calamitous.} The policy failure had calamitous consequences for the country's food supply. \emph{Calamitous consequences are extremely damaging results.}
\end{enumerate}

\vfill
\begin{center}
\small\color{Muted} Created and compiled by \textbf{Amir Kooshky}\\
\href{https://t.me/KooshkyTOEFL}{Telegram Channel}\quad\textbullet\quad
\href{https://www.instagram.com/kooshkytoefl}{Instagram}\quad\textbullet\quad
\href{https://t.me/KooshkyTOEFL\_pv}{Personal Telegram}
\end{center}
\end{document}
